\section{Introduction}

%The Paris Agreement \cite{agreement2015paris} and the global transition toward lower-carbon energy systems have reshaped investment strategies within the oil and gas industry. 
The increasing pressure to improve efficiency in hydrocarbon production has reshaped operational strategies within the oil and gas industry. While non-conventional resources such as shale oil and shale gas have gained increasing prominence, mature conventional fields continue to play a significant role in energy supply. In these environments, the main challenge is not expanding production capacity,y but improving operational efficiency under limited processing and intervention resources \cite{SIRCAR2021379}. Mathematical programming approaches have been widely applied in oil production optimization, including mixed-integer and nonlinear formulations for production allocation and artificial lift systems  \cite{doi:10.1021/ie034171z}. 
%Advances in oil and gas exploration and production technologies have enabled the large-scale development of non-conventional resources, such as shale oil and shale gas, which in many regions have gained prominence over conventional production methods. As a result, non-conventional fields have attracted increasing investment and technological innovation, reshaping production strategies within the industry. As investment and technological development increasingly concentrate on non-conventional assets, mature conventional fields face the challenge of remaining economically viable by focusing on operational efficiency and cost optimization. New technologies applied to the exploration and extraction of oil and gas have allowed the development of fracking that has overpassed conventional methods. In Argentina, oil companies are leaving conventional wells and moving to Vaca Muerta, a region that is among the five largest shale oil and gas fields in the world. %In mature conventional fields, the main challenge does not lie in expanding production, but rather in adjusting operating costs and optimizing existing resources \cite{SIRCAR2021379}. 
%A production battery in the oil industry is a set of equipment and facilities used in the initial stages of crude oA production battery consists of facilities located near the wells, where extracted fluids are separated into oil, gas, water, and other components before being transported to treatment plants.l processing at an oilfield, within the \textit{upstream} sector. 

In mature conventional oilfields, production is typically organized through multiple batteries that collect and process the output of geographically dispersed wells before delivering it to a treatment plant. A production battery consists of facilities located near the wells, where extracted fluids are separated into oil, gas, water, and other components before being transported to treatment plants. Subsequently, oil and gas are transported to a treatment plant, distribution centers, or refineries through pipelines or tanker trucks. Figure \ref{fig:deploy} shows schematically the hierarchy of a conventional oil field production elements.

\begin{figure}
    \centering
    \input{jerarquia}
    \caption{Oil field hierarchy}
    \label{fig:deploy}
\end{figure}

%Although the commercial interest is oriented to fracking wells, there are still large mature conventional fields producing oil and gas. %A production battery in the oil industry is a set of equipment and facilities used in the initial stages of crude oil processing at an oilfield, what is known as \textit{upstream} sector. These are closed to where  wells are drilled and receive the extracted product for its processing and separation of oil, natural gas, water and other elements not useful. Both oil and gas are then stored for its transportation through pipelines or tanker trucks towards distribution centers or refineries.
%The gross production arriving at the well battery contains a certain proportion of crude oil and other fluids such as water, gas, and impurities. This proportion can vary from well to well and is usually referred to as yield. The difference between gross and net volume is called loss and is generally injected back into the well through appropriate mechanisms to maintain its relative pressure  \cite{HASSAN2023181}. Enhanced Oil Recovery (EOR), uses advanced techniques like injecting fluids (CO2, steam or chemicals) into depleted oil fields to boost production beyond primary (natural pressure) and secondary (water/gas flood) methods, by altering oil/reservoir properties to improve flow, significantly increasing total recovery. 

The gross production arriving at a battery contains varying proportions of oil, water, gas, and impurities, which determine the yield of each well. The difference between gross and net volume constitutes a process loss, which in many cases is reinjected to maintain reservoir pressure conditions\cite{HASSAN2023181} . Production batteries have limited processing capacity, usually expressed in barrels or cubic meters per day, which imposes operational constraints that require periodic adjustments to the production regimes of the associated wells.
Batteries, for their part, have a gross processing capacity that is typically measured in units of volume they can process per unit of time, such as barrels per day (BPD) or cubic meters per day ($m^3/d$). The net production is obtained from the result of this processing. A battery can process the volume generated by dozens of wells, or add the volume already processed from another battery. Operational events such as maintenance or downstream constraints may require adjusting the gross volume processed by a battery, which in turn forces modifications in the operating regimes of the associated wells. This involves increasing or decreasing the production volume extracted from the wells, which translates into adjusting the operating regimes of the pumping systems.

Digital transformation initiatives have progressively incorporated supervisory and data acquisition platforms that enable remote monitoring of operational variables\cite{BIST2024100256}. However, these monitoring systems do not eliminate the need for physical intervention. In many conventional oilfields, adjustments in pumping regimes still require on-site intervention by as valve configurations or pumping flow rates. These field movements introduce significant logistical costs and operational delays, particularly in geographically dispersed fields with limited infrastructure. Therefore, production decisions cannot be treated independently from field logistics. Consequently, production reconfiguration decisions are intrinsically coupled with routing decisions: modifying the operational regime of a well implies a physical visit by maintenance personnel.

From a modeling perspective, this situation gives rise to a structured optimization problem in which continuous production control variables must be coordinated with discrete routing decisions. This structural coupling resembles integrated production–routing problems studied in operations research, such as the Production Routing Problem \cite{doi:10.1287/ijoc.1100.0439} and the Inventory Routing Problem \cite{Bertazzi2008} , where operational and logistical decisions are optimized simultaneously. In contrast, classical formulations focus on goods distribution and inventory management, whereas the problem addressed here involves endogenous routing decisions triggered by continuous operational regime adjustments.

%Conventional oil fields in Argentina are in south Patagonia mostly with sparse or no connectivity at all. The configuration set up of the wells should be performed by a worker that actually moves among the different wells. This movement has a cost both in money and time.  
%To meet the requirements arising from the change in gross production volume, production operations engineers must select the wells whose operating regimes need to be changed and calculate the new production parameters. This must be done taking into account the cost of modifying the operating regimes of each pumping unit. There is an optimization problem in this that has to complimentary objectives: i) the maximization of the oil production; ii) the minimization of the time or cost required to change the set up.
%The main contribution of this paper is the formalization of the problem through a mixed linear programming model that combines the optimization of well production subject to minimizing the distance traveled by the operator. This is achieved by combining linear programming and integer linear programming models. The model also incorporates a set of constraints related to the geographical location of the well site.

These interacting decisions naturally induce trade-offs between production efficiency and field intervention effort. Multiobjective vehicle routing problems have been extensively studied \cite{JOZEFOWIEZ2008293}, and structured mathematical programming approaches such as the $\varepsilon$-constraint method provide systematic mechanisms for generating Pareto-efficient solutions in multiobjective mixed-integer models \cite{MAVROTAS2009455} \cite{10.1007/978-3-319-15892-1_8}. However, despite advances in multiobjective routing and production planning, the explicit integration of production regime adjustment with maintenance routing under a unified multiobjective mixed-integer linear programming (MILP) framework remains largely unexplored in the context of conventional oilfield operations where routing is triggered by operational regime modification decisions.In this context, this paper proposes a multiobjective MILP formulation that integrates:

\begin{itemize}
    \item Production regime adjustment under battery capacity constraints;
    \item Maintenance team routing across geographically dispersed wells;
    \item Minimization of production losses derived from gross–net differences;
    \item Intervention cost modeling based on operational risk and technical priority.
\end{itemize}

The proposed model captures the structural interaction between production efficiency and field logistics within a unified optimization framework. Synthetic instances inspired by statistical characteristics observed in mature conventional oilfields are generated to evaluate the computational behavior of the formulation under controlled and reproducible conditions. The main contributions of this work are:
\begin{itemize}
    \item A unified multiobjective MILP formulation capturing the endogenous interaction between production regime adjustment and maintenance routing;
    \item A multiobjective structure balancing production efficiency and field intervention costs;
    \item A statistically grounded instance generation framework for reproducible experimentation;
    \item A computational analysis illustrating trade-offs between production losses and logistical costs.
\end{itemize}

%The main contribution of this work is the formulation of this problem through a mixed-integer linear programming model that integrates the optimization of production regimes with the planning of maintenance team routes. The model incorporates production, operational, and geographical constraints, and supports decision-making in conventional production environments that seek to maintain competitiveness through operational efficiency and digitalization.

The rest of the paper is organized as follows. Section 2 reviews related work. Section 3 presents the problem description. Section 4 details the MILP formulation. Section 5 describes the synthetic instance generation procedure. Section 6 reports computational results. Finally, Section 7 concludes the paper and outlines future research directions.